\section{Problem 2}
\label{sec:problem2}

The function \texttt{get\_velocity\_position\_vector} (\cref{app:p2-code}) performs the inverse operation. It takes $\Theta$ and $\mu$ and returns $[\vect{r}]_I$ and $[\ivect{V}]_I$ in PCI coordinates. The perifocal basis is reached from the inertial basis by three rotations in a 3--1--3 sequence (\cref{fig:p2}):
\begin{enumerate}
    \item A 3-rotation by $\Omega$ about $\vect{I}_z$ aligns $\vect{I}_x$ with the line of nodes $\vect{n}$.
    \item A 1-rotation by $i$ about $\vect{n}$ aligns $\vect{I}_z$ with $\ivect{h}$.
    \item A 3-rotation by $\omega$ about $\ivect{h}$ aligns $\vect{I}_x$ with $\vect{e}$, which points to periapsis.
\end{enumerate}

\begin{figure}[H]
\centering
% Problem 2: 3-1-3 rotation, inertial -> perifocal (drawn with Omega = 60, i = 35, omega = 50 deg)
% oblique projection: x -> (-0.9,-0.7), y -> (1.4,0), z -> (0,1.4)
\begin{tikzpicture}[>=Stealth, scale=1.15]
    \tikzset{old/.style={->,gray,thick}, new/.style={->,hwpurple,very thick},
        rot/.style={hwmagenta,thick},
        rotfront/.style={rot,-{Stealth[length=5pt,width=4.5pt,bend]},preaction={draw,white,line width=3pt,-}}}

    % ---------- (1) 3-rotation by Omega about I_z ----------
    \begin{scope}
        \draw[rot] (0.3,1.0) arc[start angle=0, end angle=180, x radius=0.3, y radius=0.1];
        \draw[old] (0,0) -- (-0.9,-0.7) node[below left,inner sep=1pt] {$\vect{I}_x$};
        \draw[old] (0,0) -- (1.4,0) node[right] {$\vect{I}_y$};
        \draw[old] (0,0) -- (0,1.4) node[above] {$\vect{I}_z$};
        \draw[new] (0,0) -- (0.762,-0.35) node[below right,inner sep=1pt] {$\vect{I}_x,\ \vect{n}$};
        \draw[new] (0,0) -- (1.479,0.606) node[right] {$\vect{I}_y$};
        \draw[rotfront] (-0.3,1.0) arc[start angle=180, end angle=330, x radius=0.3, y radius=0.1];
        \node[hwmagenta] at (0.55,1.0) {$\Omega$};
        \node at (0.25,-1.3) {(1) 3--$z$};
    \end{scope}

    % ---------- (2) 1-rotation by i about n ----------
    \begin{scope}[shift={(4.2,0)}]
        \draw[rot,shift={(-0.45,-0.35)},rotate=127.9] (-0.25,0) arc[start angle=180, end angle=360, x radius=0.25, y radius=0.09];
        \draw[old] (0,0) -- (-0.9,-0.7) node[below left,inner sep=1pt] {$\vect{I}_x,\ \vect{n}$};
        \draw[old] (0,0) -- (1.4,0) node[below right,inner sep=1pt] {$\vect{I}_y$};
        \draw[old] (0,0) -- (0,1.4) node[above right,inner sep=1pt] {$\vect{I}_z$};
        \draw[new] (0,0) -- (-0.803,1.147) node[above] {$\vect{I}_z,\ \ivect{h}$};
        \draw[new] (0,0) -- (1.147,0.803) node[right] {$\vect{I}_y$};
        \draw[rotfront,shift={(-0.45,-0.35)},rotate=127.9] (0.25,0) arc[start angle=0, end angle=150, x radius=0.25, y radius=0.09];
        \node[hwmagenta] at (-0.12,-0.55) {$i$};
        \node at (0.25,-1.3) {(2) 1--$x$};
    \end{scope}

    % ---------- (3) 3-rotation by omega about h ----------
    \begin{scope}[shift={(8.4,0)}]
        \draw[rot] (0.3,1.0) arc[start angle=0, end angle=180, x radius=0.3, y radius=0.1];
        \draw[old] (0,0) -- (-0.9,-0.7) node[below left,inner sep=1pt] {$\vect{I}_x,\ \vect{n}$};
        \draw[old] (0,0) -- (1.4,0) node[right] {$\vect{I}_y$};
        \draw[old] (0,0) -- (0,1.4) node[above] {$\vect{I}_z,\ \ivect{h}$};
        \draw[new] (0,0) -- (0.493,-0.45) node[below right,inner sep=1pt] {$\vect{I}_x,\ \vect{e}$};
        \draw[new] (0,0) -- (1.589,0.536) node[right] {$\vect{I}_y$};
        \draw[rotfront] (-0.3,1.0) arc[start angle=180, end angle=330, x radius=0.3, y radius=0.1];
        \node[hwmagenta] at (0.55,1.0) {$\omega$};
        \node at (0.25,-1.3) {(3) 3--$z$};
    \end{scope}
\end{tikzpicture}

\caption{3--1--3 rotation of the inertial basis $\{\vect{I}_x,\vect{I}_y,\vect{I}_z\}$ into the perifocal basis: rotation by $\Omega$ about $\vect{I}_z$, by $i$ about $\vect{n}$, then by $\omega$ about $\ivect{h}$. In each step the axes are gray before and purple after the rotation. Each axis keeps its name and is labeled with the vector it is aligned with.}
\label{fig:p2}
\end{figure}

The position and inertial velocity are first written in the perifocal basis,
\begin{equation}
r = \frac{a(1-e^2)}{1+e\cos\nu}, \qquad
\left[\vect{r}\right]_{\text{perifocal}} = \begin{bmatrix} r\cos\nu \\ r\sin\nu \\ 0 \end{bmatrix}, \qquad
\left[\ivect{V}\right]_{\text{perifocal}} = \sqrt{\frac{\mu}{p}}\begin{bmatrix} -\sin\nu \\ e+\cos\nu \\ 0 \end{bmatrix},
\end{equation}
and then transformed to PCI coordinates:
\begin{equation}
\setlength{\fboxsep}{8pt}
\boxed{
\begin{gathered}
\left[\vect{r}\right]_{\text{inertial}} = T^{\text{inertial}}_{\text{perifocal}}\left[\vect{r}\right]_{\text{perifocal}}, \qquad
\left[\ivect{V}\right]_{\text{inertial}} = T^{\text{inertial}}_{\text{perifocal}}\left[\ivect{V}\right]_{\text{perifocal}} \\[10pt]
T^{\text{inertial}}_{\text{perifocal}} = R_3(\Omega)\,R_1(i)\,R_3(\omega)
\end{gathered}
}
\end{equation}

\FloatBarrier
